\chapter{Prime Traces, Hall Operators, and Modular Conjugation}
\label{v4-arith:thermal:chapter}

The finite primary tower has one state of energy $n\log p$ at
occupation $n$. Its character is an Euler factor. To compare that
tower with a Koszul or Petersson pairing, one must also specify
its observables and the functional which measures them. We
construct those data, realize their creation operator on a Hall
sector, and determine the Hilbert space and modular conjugation
of the resulting state. The distinction between an occupation
vector, a projection, and an algebra element controls every
normalization in this passage.

\section{The prime shift and the first trace test}

Fix a prime integer $p$ and put $\epsilon=\log p>0$. The local
arithmetic datum is the vector space
\[
 V_p=\bigoplus_{n\geq0}\mathbb C v_n,
 \qquad v_n\leftrightarrow p^n,
\]
with its prime multiplication and division operators
\[
 Sv_n=v_{n+1},\qquad Rv_0=0,\qquad Rv_n=v_{n-1}\quad(n\geq1).
\]
The vectors $v_n$ are orthonormal. Their Hilbert completion is
$\mathscr H_p=\ell^2(\mathbb N_0)$, the space of sequences
$(a_n)$ with $\sum_n|a_n|^2<\infty$. Both $S$ and $R$ extend to
bounded operators of norm one, and $R=S^*$, since their matrix
coefficients are transposes of each other.

The occupation degree of $v_n$ is $n$, and its energy is
$n\epsilon$. On $V_p$ let $Nv_n=nv_n$ and $H=\epsilon N$.
The self-adjoint domain of $H$ is
\[
 \operatorname{Dom}H
 =\left\{\sum_{n\geq0}a_nv_n\in\mathscr H_p:
                  \sum_{n\geq0}n^2|a_n|^2<\infty\right\}.
\]
Indeed, testing an adjoint vector against every $v_n$ forces its
adjoint coefficients to be $n\epsilon a_n$, and this sequence must
be square summable. Thus the displayed domain is exactly the
adjoint domain. In particular, the energy has no unspecified
boundary condition at large occupation number.

Let $E=1-SR$. Direct evaluation gives the deciding identities
\begin{equation}
 RS=1,\qquad Ev_n=\delta_{n0}v_0,\qquad [R,S]=E.
 \label{v4-arith:thermal:eq:vacuum-commutator}
\end{equation}
The operator $E$ measures the probability of the vacuum.
At inverse temperature $\beta>0$, write $q=e^{-\beta\epsilon}=p^{-\beta}$.
Then $0<q<1$, and the diagonal operator $\rho_q=q^N$ has eigenvalues
$q^n$. Its partition trace and normalized Gibbs functional are
\begin{equation}
 Z_p(\beta)=\sum_{n\geq0}q^n=\frac1{1-q},\qquad
 \Phi_q(a)=(1-q)\sum_{n\geq0}q^n\langle v_n,av_n\rangle.
 \label{v4-arith:thermal:eq:gibbs-definition}
\end{equation}
The sum defining $\Phi_q$ converges absolutely for every bounded $a$,
because its absolute value is at most $\|a\|$. It is positive on
$a^*a$, and $\Phi_q(1)=1$.

For an algebra $A$, a trace means a linear functional $\tau$ such
that $\tau(ab)=\tau(ba)$ for all $a,b$. Equation
\eqref{v4-arith:thermal:eq:vacuum-commutator} gives
\begin{equation}
 \Phi_q(RS)=1,\qquad \Phi_q(SR)=q,
 \qquad \Phi_q(E)=1-q>0.
 \label{v4-arith:thermal:eq:noncyclic-test}
\end{equation}
Thus $\Phi_q$ is not a trace on the prime operator algebra. The
unnormalized thermal functional has the same obstruction:
$\operatorname{Tr}(E\rho_q)=1$.

Three uses of the word trace must consequently be kept separate.
The number $Z_p(\beta)$ is the Hilbert-space trace of an operator.
The map $a\mapsto\operatorname{Tr}(a\rho_q)$ is a thermal
functional on observables. An algebraic Hochschild trace kills
commutators. Their relation will be an explicit map with a
coefficient module, rather than an identification of these three
objects.

\section{The full algebra of prime shifts}

To determine exactly what the vacuum test excludes, we compute the
algebra generated by $S$ and $R$. All algebras and tensor products
in the algebraic calculations are over $\mathbb C$, and all sums
are finite unless convergence is specified. The operators have
cohomological degree zero. Their separate occupation degrees are
$\deg S=1$ and $\deg R=-1$.

\begin{proposition}[normal form and matrix units]
\label{v4-arith:thermal:prop:normal-form}
The algebra $\mathcal T$ generated by $S,R$ has presentation
\[
 \mathcal T=\mathbb C\langle S,R\rangle/(RS-1).
\]
Its vector-space basis is $S^iR^j$ for $i,j\geq0$, and its product is
\begin{equation}
 (S^iR^j)(S^kR^l)=
 \begin{cases}
 S^{i+k-j}R^l&j\leq k,\\
 S^iR^{j-k+l}&j\geq k.
 \end{cases}
 \label{v4-arith:thermal:eq:normal-product}
\end{equation}
The elements $E_{ij}=S^iER^j$ act as the matrix units
$E_{ij}v_n=\delta_{jn}v_i$. Their finite linear span
$\mathcal I_{\mathrm{fin}}$ is a two-sided ideal, and
\begin{equation}
 \mathcal T/\mathcal I_{\mathrm{fin}}\cong\mathbb C[z,z^{-1}],
 \qquad S\mapsto z,\quad R\mapsto z^{-1}.
 \label{v4-arith:thermal:eq:symbol-quotient}
\end{equation}
\end{proposition}

\begin{proof}
Deleting every occurrence of $RS$ reduces a word to $S^iR^j$ and
gives \eqref{v4-arith:thermal:eq:normal-product}. For linear independence, group a
finite relation by $d=i-j$. A term with this degree acts on $v_n$
as $v_{n+d}$ when $n\geq j$, and as zero otherwise. Terms of
different degrees have different output basis vectors. For a fixed
degree, successive differences between the coefficients at $v_n$
and $v_{n-1}$ recover the coefficient of the term whose threshold
is $j=n$. This proves that every coefficient in the relation is zero.

The matrix-unit formula follows from the vacuum formula. Multiplying
$E_{ij}$ on either side by $S$ or $R$ either shifts an index or gives
zero, so their span is an ideal. Finally,
\[
 E_{ij}=S^iR^j-S^{i+1}R^{j+1}.
\]
Modulo that span all normal words reduce to powers of $S$ or of
$R$, and $SR=RS=1$. The Laurent monomials are independent, so
there is no further relation and the displayed quotient is exact.
\end{proof}

The ideal contains the matrices with finite support in the specified
basis. It is not being identified with every finite-rank bounded
operator, whose vectors can have infinite support.

An occupation-homogeneous operator $a$ of degree $d$ satisfies
$[H,a]=d\epsilon a$ on $V_p$. Consequently
\[
 \alpha_t(a)=e^{itH}ae^{-itH}=e^{itd\epsilon}a.
\]
This defines a one-parameter group of automorphisms of $\mathcal T$.
Each finite sum in $\mathcal T$ has an entire continuation in $t$.
Its imaginary-time value is the automorphism
\begin{equation}
 \sigma_q(S)=qS,\qquad \sigma_q(R)=q^{-1}R,
 \qquad \sigma_q(a)=\alpha_{i\beta}(a).
 \label{v4-arith:thermal:eq:twist}
\end{equation}
It is defined on the algebraic core. No bounded extension of
$\sigma_q$ to all bounded operators is required.

\begin{proposition}[the thermal values and twisted trace identity]
\label{v4-arith:thermal:prop:thermal-values}
For $i,j\geq0$,
\begin{equation}
 \Phi_q(S^iR^j)=\delta_{ij}q^i,
 \qquad
 \Phi_q(E_{ij})=\delta_{ij}(1-q)q^i.
 \label{v4-arith:thermal:eq:thermal-normal-values}
\end{equation}
For all $a,b\in\mathcal T$,
\begin{equation}
 \Phi_q(ab)=\Phi_q\bigl(b\sigma_q(a)\bigr).
 \label{v4-arith:thermal:eq:twisted-trace}
\end{equation}
\end{proposition}

\begin{proof}
An off-diagonal normal word has zero diagonal matrix coefficients.
The word $S^iR^i$ is the projection onto the span of $v_n$ for
$n\geq i$. Summing the geometric tail proves the first formula;
subtracting adjacent tails proves the second.

It remains to verify \eqref{v4-arith:thermal:eq:twisted-trace} on basis words
$a=S^iR^j$ and $b=S^kR^l$. If $i+k\ne j+l$, both products have
nonzero occupation degree and both values vanish. Otherwise, if
$j\leq k$, then $l=i+k-j\geq i$. The two products are
$ab=S^lR^l$ and $ba=S^kR^k$, whose values satisfy
$q^l=q^{i-j}q^k$. If $j\geq k$, then $l\leq i$, and they are
$ab=S^iR^i$ and $ba=S^jR^j$, whose values satisfy
$q^i=q^{i-j}q^j$. The factor $q^{i-j}$ is exactly the action of
$\sigma_q$ on $a$.
\end{proof}

Equation \eqref{v4-arith:thermal:eq:twisted-trace} is the KMS identity on the analytic
prime core, with time evolution $\alpha_t$ and inverse temperature
$\beta$. It arises here from the actual shift multiplication and
the geometric probability distribution. In particular, a formula
assigning $p^{-\beta}$ to $RS=1$ cannot be a normalized state.
The factor $q$ is the expectation of the range projection $SR$.

\section{Arithmetic phase operators}

The shift already knows the prime through its energy. The local
arithmetic representation also contains phases. We retain them to
test whether a comparison loses more than its partition function.

Let $\Gamma_p=\mathbb Z[1/p]/\mathbb Z$, an additive group. Fix
the exponential character $\gamma\mapsto\exp(2\pi i\gamma)$;
the same character will be used when more primes are introduced. For
$\gamma\in\Gamma_p$, define
\[
 D_\gamma v_n=\exp(2\pi i p^n\gamma)v_n.
\]
The expression is independent of the rational representative of
$\gamma$. If $p^h\gamma=0$ in $\Gamma_p$, then
\begin{equation}
 D_\gamma=1+\sum_{n=0}^{h-1}
       \bigl(\exp(2\pi i p^n\gamma)-1\bigr)E_{nn}\in\mathcal T.
 \label{v4-arith:thermal:eq:phase-finite}
\end{equation}
Thus no new infinite series is needed for these operators.

\begin{proposition}[local phase relations and expectations]
\label{v4-arith:thermal:prop:phase}
The operators $D_\gamma$ are unitary and have occupation degree zero.
They satisfy
\begin{align}
 D_\gamma D_\eta&=D_{\gamma+\eta},&
 RD_\gamma S&=D_{p\gamma},\\
 SD_\gamma R&=\frac1p\sum_{p\eta=\gamma}D_\eta.&
 \label{v4-arith:thermal:eq:phase-relations}
\end{align}
Their Gibbs expectations are
\begin{equation}
 \Phi_q(D_\gamma)
  =(1-q)\sum_{n=0}^{h-1}q^n\exp(2\pi i p^n\gamma)+q^h.
 \label{v4-arith:thermal:eq:phase-thermal}
\end{equation}
\end{proposition}

\begin{proof}
The group law and adjoint formula follow from the diagonal entries.
The second relation follows on $v_n$ by replacing $p^n$ with
$p^{n+1}$. For the third, the $p$ preimages of $\gamma$ are
$\eta_0+j/p$ for $0\leq j<p$. Their phases on $v_0$ sum to zero,
because $\sum_{j=0}^{p-1}e^{2\pi i j/p}=0$. On $v_n$ for $n\geq1$,
all $p$ phases agree and equal $\exp(2\pi i p^{n-1}\gamma)$,
which is the entry of $SD_\gamma R$. The thermal formula is the
finite initial sum followed by its geometric tail.
\end{proof}

For $p=2$ and $\gamma=1/2$, these formulas become
\[
 D_{1/2}=1-2E,
 \qquad \Phi_q(D_{1/2})=2q-1.
\]
At $\beta=1$ this phase expectation is zero, while
$\Phi_q(E)=1/2$. This is a local prime calculation with a convergent
thermal density. It asserts neither a convergent global partition
trace at $\beta=1$ nor faithfulness of a representation of a larger
universal arithmetic algebra. The concrete represented algebra
generated by $S,R$ and all these $D_\gamma$ is exactly $\mathcal T$.

\section{Ordinary and twisted Hochschild quotients}

Degree-zero Hochschild homology is the universal recipient for
traces. We need its precise definition, because a thermal functional
has the twisted identity \eqref{v4-arith:thermal:eq:twisted-trace}.

An $A$-bimodule $M$ is a vector space with commuting left and right
actions of an algebra $A$. Its Hochschild chains are
$C_n(A,M)=M\otimes A^{\otimes n}$ for $n\geq0$. Write a simple
tensor as $m|a_1|\cdots|a_n$. Put $b=0$ on $C_0(A,M)$.
For $n\geq1$, the differential lowers homological degree by one
and is
\begin{align*}
 b(m|a_1|\cdots|a_n)
 &=ma_1|a_2|\cdots|a_n\\
 &\quad+\sum_{j=1}^{n-1}(-1)^j
       m|a_1|\cdots|a_ja_{j+1}|\cdots|a_n\\
 &\quad+(-1)^na_nm|a_1|\cdots|a_{n-1}.
\end{align*}
Each pair of merges in $b^2$ occurs twice with opposite signs.
Merges involving the first or last factor agree by the module
axioms, including commutation of the two actions. Therefore $b^2=0$.
In degree zero this gives
\[
 \mathrm{HH}_0(A,M)=M/\operatorname{span}\{ma-am:m\in M,a\in A\}.
\]
For $M=A$, write $\mathrm{HH}_0(A)=A/[A,A]_{\mathrm{vect}}$.
The subscript emphasizes a linear span, which need not be an ideal.
This quotient is not automatically an algebra.

For an automorphism $\sigma$ of $A$, let $M_\sigma$ be the underlying
vector space $A$ with left action $a\cdot m=am$ and right action
$m\cdot a=m\sigma(a)$. Define
\[
 \mathrm{HH}_0^\sigma(A)=\mathrm{HH}_0(A,M_\sigma)
 =A/\operatorname{span}\{ab-b\sigma(a):a,b\in A\}.
\]
This convention has exactly the cyclicity required by
\eqref{v4-arith:thermal:eq:twisted-trace}. Its first two differentials are
\begin{align*}
 b(m|a)&=m\sigma(a)-am,\\
 b(m|a|b)&=m\sigma(a)|b-m|ab+bm|a.
\end{align*}
Applying the first formula to the second cancels its six terms,
providing a direct sign check.

\begin{theorem}[the two prime Hochschild quotients]
\label{v4-arith:thermal:thm:hh-quotients}
For the algebra $\mathcal T$,
\begin{equation}
 \mathrm{HH}_0(\mathcal T)\cong\mathbb C[z,z^{-1}],
 \qquad
 [E_{ij}]=0,\quad [D_\gamma]=[1].
 \label{v4-arith:thermal:eq:ordinary-hh}
\end{equation}
For $0<q<1$ and $\sigma_q$ as in \eqref{v4-arith:thermal:eq:twist},
\begin{equation}
 \mathrm{HH}_0^{\sigma_q}(\mathcal T)\cong\mathbb C[1],
 \qquad [S^iR^j]=\delta_{ij}q^i[1].
 \label{v4-arith:thermal:eq:twisted-hh}
\end{equation}
The isomorphism in \eqref{v4-arith:thermal:eq:twisted-hh} sends $[a]$ to $\Phi_q(a)$.
It is an isomorphism of vector spaces, concentrated in occupation
degree zero. The Gibbs functional does not factor through
\eqref{v4-arith:thermal:eq:ordinary-hh}.
\end{theorem}

\begin{proof}
The Laurent quotient in Proposition~\ref{v4-arith:thermal:prop:normal-form} is
commutative, so every commutator lies in $\mathcal I_{\mathrm{fin}}$.
Conversely,
\[
 E_{ij}=[R,S^{i+1}R^j],
\]
so every matrix unit is a commutator. Thus
$[\mathcal T,\mathcal T]_{\mathrm{vect}}=\mathcal I_{\mathrm{fin}}$.
Equation \eqref{v4-arith:thermal:eq:phase-finite} gives the phase assertion.

In the twisted quotient, taking $a=S$ and
$b=S^{i-1}R^j$ for $i,j\geq1$ gives
\[
 [S^iR^j]=q[S^{i-1}R^{j-1}].
\]
The same twisted relation with $b=S^{d-1}$ gives
$(1-q)[S^d]=0$ for $d\geq1$. Taking $a=R$ and $b=R^{d-1}$ gives
$(1-q^{-1})[R^d]=0$. These scalar factors are nonzero, so every
off-diagonal word vanishes and each diagonal word has the value
in \eqref{v4-arith:thermal:eq:twisted-hh}. The class $[1]$ is nonzero because
Proposition~\ref{v4-arith:thermal:prop:thermal-values} gives a twisted trace with
$\Phi_q(1)=1$. Finally, \eqref{v4-arith:thermal:eq:noncyclic-test} forbids descent
of that functional to the ordinary quotient.
\end{proof}

Although the ordinary quotient happens to be an algebra in this
example, the twisted quotient does not inherit multiplication from
$\mathcal T$. Indeed, $[S]=[R]=0$ there, while $[RS]=[1]\ne0$.
Thus its defining subspace is not a two-sided ideal. Even the
one-dimensional answer cannot silently be endowed with the product
of representatives.

The temperature boundary distinguishes the two quotients. As
$q\uparrow1$, the values $\Phi_q(S^iR^j)$ tend to $\delta_{ij}$.
This limiting functional is the constant Laurent coefficient in
\eqref{v4-arith:thermal:eq:ordinary-hh}; it kills every matrix unit and every phase
difference $D_\gamma-1$. At the same time $Z_p(\beta)$ diverges.
Thus the appearance of an ordinary trace at the limit does not
give a finite-temperature Gibbs trace on the ordinary quotient.

\section{The density as an actual chain-level comparison}

There is a positive comparison with ordinary Hochschild chains, but
its coefficients must remember the thermal density. We construct
those coefficients directly.

For vectors $u,v$ in a Hilbert space, let $u\otimes v^*$ denote
the rank-one operator $w\mapsto\langle v,w\rangle u$, with inner
product linear in its second argument. Call an operator nuclear if
it has a representation
\[
 K=\sum_{j\geq0}u_j\otimes y_j^*,\qquad
 \sum_{j\geq0}\|u_j\|\,\|y_j\|<\infty.
\]
The series converges in operator norm. Denote the space of these
operators on $\mathscr H_p$ by $\mathcal I_1$. It is an
$\mathcal T$-bimodule: multiplying on either side by a bounded
operator preserves the summability condition.

\begin{lemma}[the trace on nuclear coefficients]
\label{v4-arith:thermal:lem:nuclear-trace}
For $K\in\mathcal I_1$, the diagonal sum
$\operatorname{Tr}(K)=\sum_n\langle v_n,Kv_n\rangle$ is absolutely
convergent, independent of its nuclear expansion and of the
orthonormal basis. It satisfies
\[
 \operatorname{Tr}(K)=\sum_j\langle y_j,u_j\rangle,
 \qquad
 \operatorname{Tr}(aK)=\operatorname{Tr}(Ka)\quad(a\text{ bounded}).
\]
\end{lemma}

\begin{proof}
For any orthonormal basis $(f_n)$, Cauchy--Schwarz and the
orthonormal expansion of the norm give
\[
 \sum_n|\langle f_n,u\rangle\langle v,f_n\rangle|
 \leq\|u\|\,\|v\|.
\]
Summing this estimate over the nuclear expansion permits exchange
of the two absolutely convergent sums. The sum over $n$ is
$\langle y_j,u_j\rangle$, proving independence of basis and
expansion. For a rank-one term, both $aK$ and $Ka$ have trace
$\langle v,au\rangle$. The same absolute estimate permits summing
this equality over all terms.
\end{proof}

The density $\rho_q=\sum_{n\geq0}q^nE_{nn}$ is nuclear, and
$\operatorname{Tr}(\rho_q)=(1-q)^{-1}$. On basis vectors,
\begin{equation}
 \rho_qa=\sigma_q(a)\rho_q\qquad(a\in\mathcal T).
 \label{v4-arith:thermal:eq:density-intertwining}
\end{equation}
All operators in this identity are bounded; it follows first for
$S,R$ and then for their finite products and sums.

\begin{theorem}[thermal coefficient comparison]
\label{v4-arith:thermal:thm:density-chain-map}
The map
\[
 K_q:M_{\sigma_q}\longrightarrow\mathcal I_1,
 \qquad m\longmapsto m\rho_q
\]
is an injective $\mathcal T$-bimodule map. It induces a chain map
\begin{equation}
 C_\bullet(\mathcal T,M_{\sigma_q})
 \longrightarrow C_\bullet(\mathcal T,\mathcal I_1),
 \qquad m|a_1|\cdots|a_n\longmapsto
             m\rho_q|a_1|\cdots|a_n.
 \label{v4-arith:thermal:eq:density-chain-map}
\end{equation}
The composite of its degree-zero part with
$(1-q)\operatorname{Tr}$ is $\Phi_q$. The chain map preserves
homological degree and homogeneous occupation degree. In the target,
occupation degree $d$ means the eigenspace for conjugation by
$e^{itN}$ with eigenvalue $e^{itd}$. No direct-sum decomposition
of all nuclear operators into these eigenspaces is asserted.
\end{theorem}

\begin{proof}
Left linearity is immediate. For the right action,
\[
 K_q(m\cdot a)=m\sigma_q(a)\rho_q
              =m\rho_qa=K_q(m)a
\]
by \eqref{v4-arith:thermal:eq:density-intertwining}. The range of $\rho_q$ contains
every basis vector up to a nonzero scalar, hence is dense. A
bounded $m$ vanishing on that range is zero; this proves injectivity.
Every face of the Hochschild differential uses multiplication or
a module action, so the bimodule identity proves
\eqref{v4-arith:thermal:eq:density-chain-map}. The trace kills the degree-one
boundaries by Lemma~\ref{v4-arith:thermal:lem:nuclear-trace}. Its value on $m\rho_q$
is the diagonal sum in \eqref{v4-arith:thermal:eq:gibbs-definition}.

Finally $\rho_q$ has occupation degree zero. Multiplication by it
preserves each homogeneous operator degree, and the chain map
does not change the number of tensor factors.
\end{proof}

The map $K_q$ is a coefficient-module map, not an algebra map:
$K_q(1)=\rho_q$ and $\rho_q^2\ne\rho_q$. It also does not map into
the algebraic core $\mathcal T$. A degree-zero element of
$\mathcal T$ has diagonal entries that are eventually constant,
whereas the entries $q^n$ of $\rho_q$ never become constant.
The nuclear coefficient module is therefore essential to this
comparison. Injectivity of the chain map is not a claim that it
is a quasi-isomorphism.

\section{A Hall-sector realization of the prime operators}

We now construct a Hall carrier on which the arithmetic shift and
division operators act. The relevant operation is given explicitly,
so its multiplication can be tested before any trace comparison.

For $n\geq0$, let
$\mathcal J_n=\mathbb C[x_1,\ldots,x_n]^{\mathfrak S_n}$, with
$\mathcal J_0=\mathbb C$. Here $\mathfrak S_n$ permutes the variables.
A polynomial of polynomial degree $j$ has cohomological degree
$2j$ and occupation degree $n$. For $f\in\mathcal J_a$ and
$g\in\mathcal J_b$, define
\begin{equation}
 (f*g)(x_1,\ldots,x_{a+b})
  =\sum_{\substack{I\sqcup J=\{1,\ldots,a+b\}\\|I|=a}}
     f(x_I)g(x_J).
 \label{v4-arith:thermal:eq:jordan-shuffle}
\end{equation}
The indices within $I,J$ are listed in increasing order; symmetry
makes that choice immaterial. We call this the Jordan shuffle
algebra. The Jordan quiver has one vertex and one loop; its
$n$-dimensional representations are endomorphisms of $\mathbb C^n$.
The formula above is the ordinary cohomological Hall shuffle
operation for this quiver. We work with its displayed polynomial
carrier and operation. Identification with a critical cohomology
construction, or with an arithmetic convolution category, is an
additional comparison and is not used below.

\begin{proposition}[the full shuffle algebra and its prime sector]
\label{v4-arith:thermal:prop:jordan-polynomial}
The product \eqref{v4-arith:thermal:eq:jordan-shuffle} is associative, commutative,
and unital. If $e_k=x^k\in\mathcal J_1$, then
\[
 \mathcal J=\bigoplus_{n\geq0}\mathcal J_n
   \cong\mathbb C[e_0,e_1,e_2,\ldots].
\]
Let $\eta_n=1\in\mathcal J_n$ and $w_n=n!\eta_n$.
The subalgebra of cohomological degree zero is
\begin{equation}
 \mathcal P=\bigoplus_{n\geq0}\mathbb C w_n,
 \qquad w_a*w_b=w_{a+b},
 \qquad \mathcal P\cong\mathbb C[x],\quad w_n\leftrightarrow x^n.
 \label{v4-arith:thermal:eq:prime-sector}
\end{equation}
\end{proposition}

\begin{proof}
Both ways to form a triple product sum once over each ordered
partition into blocks of sizes $a,b,c$. This proves associativity.
Exchanging the two blocks proves commutativity, and the empty block
gives the unit. Products of the $e_k$ are monomial symmetric
polynomials, multiplied by factorials for repeated exponents.
The monomial symmetric polynomials form a basis: every symmetric
polynomial has constant coefficient on each permutation orbit of
monomials. The factorials are nonzero and invertible in
$\mathbb C$, proving the asserted polynomial presentation.

For the constant polynomials, there are $\binom{a+b}{a}$ blocks
in \eqref{v4-arith:thermal:eq:jordan-shuffle}; thus
$\eta_a*\eta_b=\binom{a+b}{a}\eta_{a+b}$.
Multiplying by $a!b!$ proves \eqref{v4-arith:thermal:eq:prime-sector}.
\end{proof}

The factorial is forced by multiplication. For example,
$\eta_1*\eta_1=2\eta_2$, whereas $w_1*w_1=w_2$.
An identification of an orthonormal occupation basis with the
$\eta_n$ would not give the required unit shift by Hall
multiplication. This also explains the coefficient restriction:
the normalization uses $n!^{-1}$ when passing in the other direction.

Declare the $w_n$ orthonormal and complete $\mathcal P$ to
$\widehat{\mathcal P}$. Define operators on its polynomial core by
\[
 Cf=x f,\qquad
 Bf=\frac{f(x)-f(0)}x.
\]
The numerator of $Bf$ is divisible by $x$, so $B$ is a polynomial
operator. It has $Bw_0=0$ and $Bw_n=w_{n-1}$ for $n\geq1$.
Therefore $C,B$ extend boundedly, have norm one, and are adjoints.
In the unnormalized basis $B\eta_n=\eta_{n-1}/n$ for $n\geq1$.

\begin{theorem}[prime-to-Hall-module comparison]
\label{v4-arith:thermal:thm:hall-module}
The map
\begin{equation}
 U_p:V_p\longrightarrow\mathcal P,
 \qquad v_n\longmapsto w_n=n!\eta_n
 \label{v4-arith:thermal:eq:unitary-map}
\end{equation}
preserves occupation degree and cohomological degree zero. It
extends to a unitary operator
$\mathscr H_p\to\widehat{\mathcal P}$ and satisfies
\[
 U_pSU_p^{-1}=C,\qquad U_pRU_p^{-1}=B,
 \qquad U_pHU_p^{-1}w_n=n\log p\,w_n.
\]
Conjugation by $U_p$ gives a faithful algebra homomorphism
\[
 \mathcal T\longrightarrow\operatorname{End}_{\mathbb C}(\mathcal P),
 \qquad S\mapsto C,\quad R\mapsto B.
\]
It preserves the operator occupation grading and all thermal
functionals $a\mapsto\operatorname{Tr}(a e^{-\beta H})$ for
$\beta>0$. Each phase operator maps to
\[
 \widehat D_\gamma w_n=\exp(2\pi i p^n\gamma)w_n.
\]
The resulting ordinary and twisted Hochschild chain comparisons
are isomorphisms between $\mathcal T$ and its represented image.
\end{theorem}

\begin{proof}
The two bases in \eqref{v4-arith:thermal:eq:unitary-map} are orthonormal and
indexed by the same nonnegative integers. Evaluation on each
basis vector proves all three intertwining equations, including
equality of the maximal energy domains. Conjugation preserves
products and is injective. It also sends $E_{ij}$ to the operator
taking $w_j$ to $w_i$ and annihilating the other basis vectors,
so the finite phase formula is preserved.

The thermal trace is an absolutely convergent diagonal sum, and
the two diagonal sums have identical terms under $U_p$. The same
holds for $H^r e^{-\beta H}$ for each fixed nonnegative integer
$r$, since $\sum_n n^r q^n$ converges. Finally, applying the
conjugation map to every tensor factor commutes with every
Hochschild face. It intertwines $\sigma_q$ by the occupation
grading, giving the twisted-chain assertion as well.
\end{proof}

This comparison includes the adjoint prime operation and the
arithmetic phases, beyond a character identity. Its target is an
algebra of operators on a Hall sector. The following obstruction
prevents replacing that target by Hall multiplication itself.

\begin{theorem}[failure of a trace-preserving Hall-algebra map]
\label{v4-arith:thermal:thm:hall-algebra-obstruction}
Let $A$ be a commutative unital algebra. Every unital algebra map
$f:\mathcal T\to A$ kills $E$ and every $E_{ij}$. For $0<q<1$,
there is no linear functional $\lambda:A\to\mathbb C$ with
$\lambda\circ f=\Phi_q$. In particular, the prime-to-Hall-module
comparison cannot be replaced by a Gibbs-preserving algebra map
from $\mathcal T$ into $\mathcal J$.

There is also no unital algebra map $\mathcal T\to\mathcal J$
that sends $S$ to $e_0$. The algebra map
$\mathbb C[S]\to\mathcal P$, $S\mapsto e_0$, does exist, but
does not extend across the relation $RS=1$.
\end{theorem}

\begin{proof}
The image of $[R,S]$ is zero in a commutative algebra, so
\eqref{v4-arith:thermal:eq:vacuum-commutator} gives $f(E)=0$. Multiplication on
both sides then kills all matrix units. But
$\Phi_q(E)=1-q\ne0$, which excludes the functional $\lambda$.
For the second assertion, $f(R)e_0=1$ would make the polynomial
variable $e_0$ invertible. Evaluation at $e_0=0$ and all other
variables zero gives $0=1$, a contradiction.
\end{proof}

The operator $B$ is neither multiplication by a polynomial nor a
derivation. It kills $1$ but has $B(x)=1$, excluding a multiplication
operator. Also $B(x^2)=x$, while a derivation with $B(x)=1$ would
give $2x$. It is the backward shift required by prime division.
This distinction is precisely where the operator realization
contains more structure than the commutative Hall sector.

\section{A finite native Iwahori calculation}

The prime shift is explicit arithmetic operator data. A local
Iwahori presentation supplies a different, finite test which must
also be respected by any proposed native comparison.

Let $X_p$ be the set of one-dimensional subspaces of $\mathbb F_p^2$.
There are $p+1$ such lines: $p$ of the form $\mathbb F_p(1,a)$
and the line $\mathbb F_p(0,1)$. The group
$G_p=\mathrm{GL}_2(\mathbb F_p)$ acts transitively on ordered pairs
of distinct lines, since two distinct lines admit basis vectors.
Let $W_p=\mathbb C^{X_p}$, with its standard orthonormal basis,
and let $T$ have matrix entry one between distinct lines and
zero on the diagonal.

\begin{proposition}[the finite Hecke algebra and its trace]
\label{v4-arith:thermal:prop:finite-hecke}
The algebra of operators on $W_p$ commuting with $G_p$ is
\begin{equation}
 \mathcal A_p=\mathbb C1\oplus\mathbb CT,
 \qquad T^2=(p-1)T+p1.
 \label{v4-arith:thermal:eq:hecke-relation}
\end{equation}
It is the finite simple-reflection Hecke algebra with parameter $p$.
Its primitive idempotents, its degree-zero Hochschild homology,
and its normalized matrix trace are
\begin{align}
 e_+&=\frac{T+1}{p+1},&e_-&=\frac{p-T}{p+1},
       &\mathcal A_p&=\mathbb Ce_+\oplus\mathbb Ce_-,\\
 \mathrm{HH}_0(\mathcal A_p)&=\mathcal A_p,
       &\tau_p(a)&=\frac{\operatorname{Tr}_{W_p}(a)}{p+1},
       &\tau_p(e_+)&=\frac1{p+1},\quad
         \tau_p(e_-)=\frac p{p+1}.
 \label{v4-arith:thermal:eq:hecke-trace}
\end{align}
For every integer $r\geq0$,
$\tau_p(T^r)=(p^r+p(-1)^r)/(p+1)$.
\end{proposition}

\begin{proof}
A commuting matrix has coefficients constant on the two $G_p$-orbits
in $X_p\times X_p$: equal and distinct pairs. Thus it is a linear
combination of $1$ and $T$. For a diagonal entry of $T^2$ there
are $p$ possible intermediate lines; for an off-diagonal entry
there are $p-1$. This proves \eqref{v4-arith:thermal:eq:hecke-relation} and gives
the entire multiplication table.

On the constant vectors $T$ acts by $p$, and on the sum-zero
subspace it acts by $-1$. These subspaces have dimensions $1$
and $p$, respectively. The formulas for $e_+,e_-$ are their
projections; they are nonzero, orthogonal, sum to one, and span
the algebra. The trace and power formulas follow from the two
eigenvalues with these multiplicities. The algebra is commutative,
so its degree-zero Hochschild quotient is itself.
\end{proof}

This calculation is tied to the local Iwahori relation without a
categorical assumption. Here
$\mathbb Z_p=\varprojlim_h\mathbb Z/p^h\mathbb Z$ is the ring
of compatible residue sequences and $\mathbb Q_p$ is its fraction
field. If $K=\mathrm{GL}_2(\mathbb Z_p)$ and
$I\subset K$ is the inverse image of the upper triangular subgroup
under reduction modulo $p$, then $K/I$ is $X_p$.
The $K$-invariant kernels on $(K/I)^2$ are exactly the two kernels
above. Composition sums over the intermediate coset and gives
\eqref{v4-arith:thermal:eq:hecke-relation}. More explicitly, on functions constant
on left and right $I$-cosets and supported on finitely many such
cosets, the convolution is
\[
 (f*g)(x)=\sum_{yI\in\mathrm{GL}_2(\mathbb Q_p)/I}
                    f(y)g(y^{-1}x).
\]
Only finitely many summands are nonzero, and the summand is
independent of the chosen representative $y$. This is the
normalization in which the characteristic function of $I$ is the
unit, or equivalently $\operatorname{vol}(I)=1$. Extension by zero
from $K$ preserves this multiplication, since a product of two
elements of $K$ lies in $K$. Thus
$\mathcal A_p$ is an actual finite subalgebra of the local
Iwahori function algebra. No sheaf-to-function or global
categorical trace equivalence is used.

\begin{theorem}[finite Hecke obstructions to the polynomial target]
\label{v4-arith:thermal:thm:hecke-obstruction}
There is no unital algebra map $f:\mathcal A_p\to\mathcal J$
and linear functional $\lambda:\mathcal J\to\mathbb C$ such that
$\lambda\circ f=\tau_p$.
More generally, a unital map from $\mathcal A_p$ to any algebra
with no idempotents except zero and one cannot preserve $\tau_p$
through a linear functional.

There is no grading on $\mathcal A_p$ by nonnegative integers
with $1$ in degree zero and $T$ homogeneous of positive degree.
In particular, reflection length does not supply an additive
positive occupation energy for this Hecke multiplication.
\end{theorem}

\begin{proof}
The two idempotents have nonzero trace and sum to one. In an
algebra with only trivial idempotents, one of their images is
zero. A linear functional then vanishes on that image, contradicting
\eqref{v4-arith:thermal:eq:hecke-trace}. The polynomial algebra $\mathcal J$ is
an integral domain: any two nonzero polynomials lie in a polynomial
ring on finitely many variables, where their leading monomials
have nonzero product. An idempotent $e$ in a domain satisfies
$e(e-1)=0$, hence is zero or one.

For the grading assertion, the degree-zero component of
\eqref{v4-arith:thermal:eq:hecke-relation} would say $0=p1$, which is impossible
over $\mathbb C$. Equivalently,
$T^{-1}=p^{-1}(T-(p-1)1)$, so $T$ is invertible and cannot be a
positive-degree element of a nonnegatively graded unital algebra.
\end{proof}

The length filtration does exist: put $F_0\mathcal A_p=\mathbb C1$
and $F_n\mathcal A_p=\mathcal A_p$ for $n\geq1$. Its associated
graded algebra is
$\mathbb C[t]/(t^2)$ with $\deg t=1$. This is a change of
multiplication. There is no filtered algebra isomorphism recovering
$\mathcal A_p$ from this graded algebra, since the latter has a
nonzero nilpotent and the former is a product of fields.

There is an elementary comparison with simple-object cohomology.
A representation of the quiver with two isolated vertices is a
pair of vector spaces; its subrepresentations are pairs of
subspaces. It is simple precisely when one space is one-dimensional
and the other is zero. Thus its two simple isomorphism classes,
of dimension vectors $(1,0)$ and $(0,1)$, form two points. The
cohomology of each point is $\mathbb C$ in degree zero and zero
in other degrees. Their direct sum is
$B=\mathbb C b_+\oplus\mathbb C b_-$. The linear map
\[
 e_+\mapsto b_+,\qquad e_-\mapsto b_-
\]
is an isomorphism from $\mathrm{HH}_0(\mathcal A_p)$ to $B$;
the functional with values $(p+1)^{-1}$ and $p(p+1)^{-1}$ preserves
$\tau_p$. It is a comparison of cohomological-degree-zero vector
spaces. The two dimension-vector labels are additional labels on
its target, not a multiplicative grading on $\mathcal A_p$.
Indeed $e_+^2=e_+$ already contradicts additive positive degree.
Nor does this linear map supply an energy on the native Hecke
trace. Identifying this simple-object cohomology with BPS
cohomology, with a specified potential, shifts, and critical
functor, remains a separate construction; it does not follow
from the two-point calculation.

\section{Tensor shift algebras over finite prime sets}

The operator comparison is compatible with adding primes. For a
finite set $F$, let
\[
 V_F=\bigotimes_{p\in F}V_p,\qquad
 \mathcal T_F=\bigotimes_{p\in F}\mathcal T_p,\qquad
 \mathcal P_F=\bigotimes_{p\in F}\mathcal P_p.
\]
These are algebraic tensor products, with occupation multi-degree
$\mathbf n=(n_p)$. Complete the basis vectors to Hilbert spaces
and give them energy $\sum_{p\in F}n_p\log p$. The tensor product
of the maps $U_p$ is unitary, intertwines each shift, adjoint and
phase, and preserves energy. All these assertions follow on the
tensor basis; no exchange of infinite products is involved.
The thermal density and normalized state are
\[
 \rho_{F,\beta}=\bigotimes_{p\in F}p^{-\beta N_p},\qquad
 Z_F(\beta)=\prod_{p\in F}(1-p^{-\beta})^{-1},\qquad
 \Phi_{F,\beta}=\bigotimes_{p\in F}\Phi_{p^{-\beta}}.
\]
The finite product is convergent for every $\beta>0$.

\begin{proposition}[finite-prime Hochschild comparisons]
\label{v4-arith:thermal:prop:finite-prime}
For the componentwise imaginary-time automorphism $\sigma$,
\[
 \mathrm{HH}_0(\mathcal T_F)
     \cong\mathbb C[z_p^{\pm1}:p\in F],
 \qquad
 \mathrm{HH}_0^\sigma(\mathcal T_F)\cong\mathbb C.
\]
The second isomorphism is the normalized product Gibbs functional.
For $F\subset G$, the algebra inclusion
$a\mapsto a\otimes1$ induces the Laurent-polynomial inclusion
on ordinary quotients and the identity on the normalized twisted
quotients. For $F=\varnothing$, all algebras and quotients in
these formulas are $\mathbb C$ and the partition trace is one.
\end{proposition}

\begin{proof}
For unital algebras $A,B$, their commutator subspace in $A\otimes B$
is
\[
 [A,A]\otimes B+A\otimes[B,B].
\]
The inclusion from right to left follows by commuting with a factor
of the form $a\otimes1$ or $1\otimes b$. For the reverse inclusion,
expand a simple-tensor commutator as
\[
 aa'\otimes bb'-a'a\otimes b'b
   =[a,a']\otimes bb'+a'a\otimes[b,b'].
\]
Taking vector-space quotients proves the ordinary formula by
induction.

In the twisted quotient, the relations in each factor may be used
while holding all other factors fixed. The reductions in the proof
of Theorem~\ref{v4-arith:thermal:thm:hh-quotients} therefore kill a normal tensor
word unless each factor is diagonal; in that case its class is
$\prod_p p^{-\beta i_p}[1]$. The product Gibbs functional gives
the matching nonzero class. It takes value one on each new unit,
which proves compatibility under $F\subset G$. The empty case
is the usual empty tensor product.
\end{proof}

\section{Arithmetic phases across two primes}

Tensoring the shift algebras preserves their local phase operators,
but those operators must be distinguished from the phase of the
integer represented by all occupation numbers together. This
provides a concrete test of the prime-set comparison itself.

For a finite set $F$ of primes, let
\[
 \Gamma_F=\mathbb Z[p^{-1}:p\in F]/\mathbb Z,
 \qquad M(\mathbf n)=\prod_{p\in F}p^{n_p}.
\]
On the occupation basis define the arithmetic phase
\[
 D^{F}_\gamma v_{\mathbf n}
    =\exp(2\pi i M(\mathbf n)\gamma)v_{\mathbf n},
       \qquad \gamma\in\Gamma_F.
\]
Let $\mathcal A_F^{\mathrm{arith}}$ be the concrete algebra of
operators generated by all $S_p,R_p$ and these $D^F_\gamma$.
Each generator preserves the algebraic occupation space and is
bounded on its Hilbert completion. The phases have occupation
multi-degree zero. For a singleton $F=\{p\}$ this algebra is
$\mathcal T_p$, by \eqref{v4-arith:thermal:eq:phase-finite}.

For every $p\in F$, direct evaluation on the integer $M(\mathbf n)$
gives
\[
 R_pD^F_\gamma S_p=D^F_{p\gamma},\qquad
 S_pD^F_\gamma R_p=\frac1p\sum_{p\eta=\gamma}D^F_\eta.
\]
In the second identity the root sum is zero when $p$ does not
divide $M(\mathbf n)$, and otherwise is the phase of
$M(\mathbf n)/p$. Thus these are actual arithmetic phase
relations on the smooth-integer basis, rather than relations
imposed on a tensor product of independent phases.

\begin{proposition}[the two-prime phase obstruction]
\label{v4-arith:thermal:prop:two-prime-phase}
Put $F=\{3,2\}$ and $\zeta=\exp(2\pi i/3)$. The operator
$D^F_{1/3}$ does not belong to $\mathcal T_3\otimes\mathcal T_2$.
There is no algebra map
\[
 \iota:\mathcal A_{\{3\}}^{\mathrm{arith}}
       \longrightarrow\mathcal A_{\{3,2\}}^{\mathrm{arith}}
\]
which simultaneously sends $S_3,R_3$ to their same-prime operators
and $D^{\{3\}}_{1/3}$ to $D^{\{3,2\}}_{1/3}$.
Writing $a=3^{-\beta}$ and $b=2^{-\beta}$, its normalized thermal
expectation for $\beta>0$ is
\begin{equation}
 \Phi_{\{3,2\},\beta}(D^F_{1/3})
   =a+(1-a)\frac{\zeta+b\zeta^2}{1+b}.
 \label{v4-arith:thermal:eq:two-prime-phase-trace}
\end{equation}
Here $\Phi$ denotes the Gibbs functional of the same diagonal
finite-prime density, now evaluated on the enlarged operator algebra.
\end{proposition}

\begin{proof}
Write the occupation indices as $(n,m)$ for $3^n2^m$. If $n\geq1$,
the phase is one. If $n=0$, then $2^m$ is congruent to $1$ or $2$
modulo $3$ according to the parity of $m$. The phase entries are
therefore $\zeta,\zeta^2,\zeta,\zeta^2,\ldots$ on that row.

For any element of $\mathcal T_3\otimes\mathcal T_2$, its diagonal
matrix coefficients on the row $(0,m)$ are eventually constant.
To see this, use its finite expansion into normal tensor words.
A diagonal contribution requires equal creation and division
indices in each factor. In the first factor only index zero
survives at occupation zero. In the second, $S_2^iR_2^i$ has
diagonal entry $1$ for $m\geq i$ and $0$ before that threshold.
A finite sum of these sequences is eventually constant.
The alternating phase sequence is not, proving nonmembership.

In the singleton algebra, the exact identity is
\[
 D^{\{3\}}_{1/3}=1+(\zeta-1)(1-S_3R_3).
\]
Any map preserving $S_3,R_3$ must preserve this identity. Its
right side has value $\zeta$ on every vector $(0,m)$, whereas
the proposed image on the left has value $\zeta^2$ when $m=1$.
Since $\zeta\ne\zeta^2$, the stated map cannot exist.

The probability that $n\geq1$ is $a$. Conditional on $n=0$,
the normalized sum over the other occupation number is
\[
 (1-b)\sum_{m\geq0}b^m\zeta^{2^m}
    =(1-b)\frac{\zeta+b\zeta^2}{1-b^2}
    =\frac{\zeta+b\zeta^2}{1+b}.
\]
All sums converge absolutely, which proves
\eqref{v4-arith:thermal:eq:two-prime-phase-trace}. For comparison, the singleton
expectation is $a+(1-a)\zeta$.
\end{proof}

At $\beta=1$, the singleton expectation equals $i\sqrt3/3$,
whereas the two-prime expectation is $i\sqrt3/9$. Thus even
normalized phase expectations do not obey the independent-factor
rule. The compatible states of Proposition~\ref{v4-arith:thermal:prop:finite-prime}
concern the tensor shift algebras and their factorwise phases;
they are not a compatible phase-preserving system of the enlarged
arithmetic algebras just defined.

There is still a positive operator comparison. The unitary
$U_F=\bigotimes_{p\in F}U_p$ sends $D^F_\gamma$ to the operator
on $\mathcal P_F$ which multiplies the monomial $w_{\mathbf n}$ by
$\exp(2\pi i M(\mathbf n)\gamma)$. Together with the creation
and backward-shift operators, this gives a faithful representation
of $\mathcal A_F^{\mathrm{arith}}$ on the Hall sector. It preserves
occupation degrees, energy, and all finite-temperature thermal
operator traces, by the same basis calculation as in
Theorem~\ref{v4-arith:thermal:thm:hall-module}.

The density comparison also survives this enlargement. Imaginary
time scales each $S_p,R_p$ by $p^{-\beta},p^{\beta}$ and fixes
every phase. Every algebra element is a finite sum of homogeneous
bounded operators, so this defines an algebra automorphism $\sigma$.
On every occupation vector,
$\rho_{F,\beta}c=\sigma(c)\rho_{F,\beta}$ for
$c\in\mathcal A_F^{\mathrm{arith}}$. The map
$m\mapsto m\rho_{F,\beta}$ is therefore an injective bimodule map
from the right-twisted coefficient module to nuclear operators,
and gives the Hochschild chain map already proved in
Theorem~\ref{v4-arith:thermal:thm:density-chain-map}. This is not a computation of
the enlarged algebra's ordinary or twisted Hochschild homology.
The formulas of Proposition~\ref{v4-arith:thermal:prop:finite-prime} retain their
stated carrier $\mathcal T_F$.

\section{The prime Gibbs Hilbert space and its modular operators}
\label{v4-arith:thermal:gns-section}

The occupation Hilbert space carries the observables. The Hilbert
space of a state is a different object: it completes observables
in the norm measured by that state. Computing this second space
will determine the modular half-density and the energy available
for a pairing comparison.

Fix $p$, $\beta>0$, $q=p^{-\beta}$, and the normalized density
\[
 \varrho_q=(1-q)\rho_q,\qquad d_j=(1-q)q^j.
\]
Thus $\Phi_q(x)=\operatorname{Tr}(x\varrho_q)$ and
$\sum_jd_j=1$. On $\mathcal T$ use the inner product linear in
its second argument,
\begin{equation}
 \langle[x],[y]\rangle_q=\Phi_q(x^*y),\qquad
 \|[x]\|_q^2=\sum_{j\geq0}d_j\|xv_j\|^2.
 \label{v4-arith:thermal:gns-norm}
\end{equation}
Every $d_j$ is strictly positive. Consequently the norm is zero
only for $x=0$. Let $\mathscr G_q$ be its Hilbert completion,
let $\Omega_q=[1]$, and let $\pi_q(a)[x]=[ax]$.
Since $\|[ax]\|_q\leq\|a\|\|[x]\|_q$, this action extends
boundedly to $\mathscr G_q$. The vector $\Omega_q$ is cyclic
by definition. These data are the Gelfand--Naimark--Segal, or GNS,
representation of the specified state.

A Hilbert--Schmidt operator on $\mathscr H_p$ is a matrix $X$
with $\sum_{i,j}|X_{ij}|^2<\infty$. Such a matrix defines a
bounded operator: Cauchy--Schwarz gives
$\|Xv\|\leq(\sum_{i,j}|X_{ij}|^2)^{1/2}\|v\|$ first on
finite-support vectors and then by continuity. The space
$\mathcal S_2(\mathscr H_p)$ has inner product
$\langle X,Y\rangle_2=\sum_{i,j}\overline{X_{ij}}Y_{ij}$.
The matrix units $E_{ij}$ form its orthonormal basis.

\begin{theorem}[explicit GNS carrier]
\label{v4-arith:thermal:gns-carrier}
There is a unitary map
\begin{equation}
 \mathcal W_q:\mathscr G_q\longrightarrow\mathcal S_2(\mathscr H_p),
 \qquad [x]\longmapsto x\varrho_q^{1/2}.
 \label{v4-arith:thermal:gns-unitary}
\end{equation}
It identifies $\pi_q(a)$ with left multiplication by $a$ and
$\Omega_q$ with the diagonal matrix
$\varrho_q^{1/2}=\sum_j\sqrt{d_j}E_{jj}$. In particular,
\[
 \langle[E_{ij}],[E_{kl}]\rangle_q
       =\delta_{ik}\delta_{jl}d_j,
 \qquad
 \langle[S^r],[S^s]\rangle_q=\delta_{rs}
       \quad(r,s\geq0).
\]
The strong operator closure of $\pi_q(\mathcal T)$ is left
multiplication by $B(\mathscr H_p)$. Its commutant consists of
right multiplications by bounded operators, and $\Omega_q$ is
separating for the left algebra.
\end{theorem}

\begin{proof}
The matrix of $x\varrho_q^{1/2}$ has column $j$ equal to
$\sqrt{d_j}xv_j$. Its squared Hilbert--Schmidt norm is
\eqref{v4-arith:thermal:gns-norm}, so $\mathcal W_q$ is an
isometry. Since $\mathcal W_q([E_{ij}])=\sqrt{d_j}E_{ij}$,
its range contains every finite matrix and is dense. It extends
unitarily to the completions. The left action and vacuum formulas
follow directly. Matrix-unit multiplication gives the first inner
product formula. For the second, $R^rS^s$ is the identity when
$r=s$ and a nonconstant pure shift or division word otherwise;
its expectation is therefore $\delta_{rs}$.

For the closure assertion, the finite matrices lie in $\mathcal T$.
If $P_n=\sum_{j=0}^nE_{jj}$ and $a$ is bounded, then
$P_naP_n\to a$ strongly and the norms are uniformly bounded.
Left multiplication converges strongly on Hilbert--Schmidt
operators: verify it on finite matrices and approximate an
arbitrary Hilbert--Schmidt matrix in its norm. Thus the strong
closure contains every bounded left multiplication.

Identify $E_{ij}$ with $v_i\otimes\overline{v_j}$ in
$\mathscr H_p\otimes\overline{\mathscr H_p}$. Commuting with
all left diagonal matrix units forces a bounded operator to
preserve each row subspace. Commuting with the left off-diagonal
matrix units forces the same bounded column action on every
row. Hence the commutant is exactly the bounded action on the
second tensor factor, or equivalently all right multiplications.
Applying this argument with the two factors interchanged shows
that the left algebra is strongly closed, proving the asserted
equality. Finally, if $a\varrho_q^{1/2}=0$, then $av_j=0$ for
every $j$ because $d_j>0$; hence $a=0$.
\end{proof}

The state is thus represented in the standard left-right form
of the full operator algebra on a countable Hilbert space.
In particular, its local von Neumann algebra is a type~I factor:
it has scalar centre and minimal projections given by left
multiplication by $E_{jj}$. No type~III assertion about an
infinite-prime system applies to this single-prime carrier.

\begin{theorem}[modular operator, conjugation, and the half-density]
\label{v4-arith:thermal:modular}
In the Hilbert--Schmidt representation, define
\[
 J_qX=X^*,\qquad
 \Delta_q E_{ij}=q^{i-j}E_{ij},
\]
where $\Delta_q$ has its maximal diagonal self-adjoint domain
\[
 \operatorname{Dom}\Delta_q
   =\left\{X:\sum_{i,j}q^{2(i-j)}|X_{ij}|^2<\infty\right\}.
\]
The conjugate-linear operator $[x]\mapsto[x^*]$ on
$\pi_q(\mathcal T)\Omega_q$ is closable, and its closure is
$\mathfrak S_q=J_q\Delta_q^{1/2}$, with domain
$\sum_{i,j}q^{i-j}|X_{ij}|^2<\infty$.
The map $J_q$ is an antiunitary involution, fixes $\Omega_q$, and
conjugates the left operator algebra onto its commutant.
For $a\in\mathcal T$ and $t\in\mathbb R$,
\begin{equation}
 \Delta_q^{it}[a]=[\alpha_{-\beta t}(a)],
 \qquad
 J_q[S^r]=q^{-r/2}[R^r],\qquad
 J_q[R^r]=q^{r/2}[S^r].
 \label{v4-arith:thermal:modular-formulas}
\end{equation}
For the phase operators, $J_q[D_\gamma]=[D_{-\gamma}]$.
\end{theorem}

\begin{proof}
Adjoint transposes and conjugates matrix coefficients, so $J_q$
is antiunitary and $J_q^2=1$. The displayed diagonal operator
is positive and self-adjoint by the same adjoint-domain test as
for $H$. Its square root has the asserted domain. On a matrix
$X=x\varrho_q^{1/2}$ from the algebraic GNS domain,
\[
 (\Delta_q^{1/2}X)_{ij}=\sqrt{d_i}\,x_{ij},
 \qquad
 (J_q\Delta_q^{1/2}X)_{ij}
       =\overline{x_{ji}}\sqrt{d_j}.
\]
The latter is the matrix of $x^*\varrho_q^{1/2}$. The former is
Hilbert--Schmidt since
$\sum_{i,j}d_i|x_{ij}|^2\leq\|x\|^2\sum_i d_i$.
Thus the initial conjugate-linear operator is a restriction of
the closed operator $J_q\Delta_q^{1/2}$. Its domain contains all
finite matrices, which form a core for this weighted diagonal
operator by truncating the convergent weighted square sum.
The closure is therefore exactly the stated operator.

The positive factor is $\Delta_q^{1/2}$, because composition
with the antiunitary $J_q$ preserves the norm and hence the
quadratic form of the positive factor. This proves the polar
decomposition directly. The diagonal vacuum is fixed by adjoint,
and $J_qL_aJ_q$ is right multiplication by $a^*$. The preceding
commutant computation proves its range assertion.

On $E_{ij}$ the modular group has scalar
$q^{it(i-j)}=e^{-it\beta(i-j)\log p}$. This is the physical
time evolution at time $-\beta t$. The creation vector $[S^r]$
has modular eigenvalue $q^r$, so
$\mathfrak S_q[S^r]=[R^r]=q^{r/2}J_q[S^r]$.
The division formula follows in the same way, or from $J_q^2=1$.
Every phase is diagonal, so its modular eigenvalue is one and
adjoint replaces $\gamma$ with $-\gamma$.
\end{proof}

At $\beta=1$ and $m=p^r$, the creation formula is
\[
 J_q(\mu_m\Omega_q)=m^{1/2}\mu_m^*\Omega_q.
\]
The sign of this exponent is forced by norms:
$\|\mu_m\Omega_q\|=1$, whereas
$\|\mu_m^*\Omega_q\|=m^{-1/2}$.
Using the inverse half-density on the creation vector would give
norm $m^{-1}$, contradicting antiunitarity for $m>1$.
The inverse half-density belongs to the division vector instead.

\begin{proposition}[occupation energy and modular energy]
\label{v4-arith:thermal:modular-energy}
Let $L_q=-\beta^{-1}\log\Delta_q$ on its maximal diagonal
self-adjoint domain. Then
\[
 L_qE_{ij}=(i-j)\log p\,E_{ij},\qquad J_qL_qJ_q=-L_q.
\]
It has both positive and negative eigenvalues and an
infinite-dimensional zero eigenspace. For every nonzero real $s$,
$e^{-sL_q}$ is unbounded and hence cannot have a nuclear partition
trace. The nonnegative occupation Hamiltonian $H$ on
$\mathscr H_p$ is therefore not the modular Hamiltonian on the
full GNS space.
\end{proposition}

\begin{proof}
Taking the real logarithm of each positive eigenvalue gives the
formula and the domain
$\sum_{i,j}(i-j)^2(\log p)^2|X_{ij}|^2<\infty$.
Adjoint exchanges $i$ and $j$. The vectors $E_{10},E_{01}$
have opposite nonzero energies, while every $E_{ii}$ has energy
zero. If $s>0$, the eigenvalues of $e^{-sL_q}$ on $E_{0j}$
are $p^{sj}$ and tend to infinity; for $s<0$ use $E_{i0}$.
The trace claim follows because nuclear operators are bounded.
\end{proof}

One weighted lattice does occur inside the GNS space, with a
specific variance. The vectors $[R^r]$, rather than $[S^r]$,
satisfy
\[
 \langle[R^r],[R^s]\rangle_q=\delta_{rs}q^r.
\]
Consequently the map from a lattice basis $f_r$ with
$\langle f_r,f_s\rangle=\delta_{rs}q^r$ to $[R^r]$ is
isometric. It reverses occupation degree and physical energy.
Multiplying $[R^r]$ by $q^{-r/2}$ gives an orthonormal family,
but does not turn it into the creation family under the same
energy grading. Any comparison with a native Koszul or Petersson
pairing must specify which family it identifies and prove its
intertwining equations. Equality of partition functions supplies
none of these maps.

For $F=\varnothing$, the GNS space is $\mathbb C\Omega$,
$J(z\Omega)=\overline z\Omega$, and $\Delta=1$. For a
finite prime set $F$, all these constructions tensor:
$\mathscr G_{F,\beta}=\bigotimes_{p\in F}\mathscr G_{p^{-\beta}}$,
$\Omega_{F,\beta}=\bigotimes_p\Omega_{p^{-\beta}}$, and
$J_{F,\beta}=\bigotimes_pJ_{p^{-\beta}}$. The matrix-unit proof
gives
\[
 \Delta_{F,\beta}E_{\mathbf i,\mathbf j}
   =\prod_{p\in F}p^{-\beta(i_p-j_p)}E_{\mathbf i,\mathbf j}.
\]
For $F\subset G$, the map
$X\mapsto X\otimes\varrho_{G\setminus F,\beta}^{1/2}$
is an isometry intertwining the vacuum, $J$, and $\Delta^{it}$.
Indeed the new vacuum has Hilbert--Schmidt norm one, is fixed
by adjoint, and has modular eigenvalue one. Thus the incomplete
Hilbert tensor product with these reference vectors carries a
well-defined antiunitary involution extending the finite $J$'s.
This construction uses the tensor shift algebras. The arithmetic
phase obstruction of \Cref{v4-arith:thermal:prop:two-prime-phase}
prevents identifying it, by the displayed local phase quotients,
with an adelic arithmetic GNS realization.

\section{A bilinear pairing and the variance of modular conjugation}
\label{v4-arith:thermal:pairing-section}

A Hermitian pairing and a bilinear Koszul pairing use conjugation
in different places. The GNS carrier has a concrete bilinear
pairing which makes this distinction visible. For
$X,Y\in\mathcal S_2(\mathscr H_p)$ put
\[
 B_q(X,Y)=\sum_{i,j}X_{ij}Y_{ji},\qquad
 h_q(X,Y)=\sum_{i,j}X_{ij}\overline{Y_{ij}}.
\]
Here $h_q$ is linear in its first variable, in the convention of
the Petersson integral; it is the reverse of the GNS inner-product
convention used above. Cauchy--Schwarz proves absolute convergence
of the first sum and $|B_q(X,Y)|\leq\|X\|_2\|Y\|_2$.

\begin{proposition}[bilinear-to-Hermitian comparison]
\label{v4-arith:thermal:bilinear-pairing}
The form $B_q$ is a continuous nondegenerate complex-bilinear form,
and
\[
 B_q(X,J_qY)=h_q(X,Y),\qquad h_q(X,X)=\|X\|_2^2.
\]
In the cyclic matrix-unit classes,
\[
 B_q([E_{ij}],[E_{kl}])
   =\delta_{jk}\delta_{il}\sqrt{d_i d_j},\qquad
 h_q([E_{ij}],[E_{kl}])=\delta_{ik}\delta_{jl}d_j.
\]
Inserting $J_q$ into the second argument of $h_q$ instead produces
a bilinear form, not a Hermitian form.
\end{proposition}

\begin{proof}
Adjoint gives $(J_qY)_{ji}=\overline{Y_{ij}}$, proving the first
identity by the absolutely convergent sum. It also proves
nondegeneracy: for $X\ne0$, the choice $Y=J_qX$ gives
$B_q(X,Y)=\|X\|_2^2>0$. The cyclic matrix-unit formulas follow
from $[E_{ij}]\mapsto\sqrt{d_j}E_{ij}$.
Finally $J_q$ is conjugate-linear, so composing it with the
conjugate-linear second slot of $h_q$ makes that slot linear.
For example, replacing $Y$ by $iY$ multiplies $h_q(X,Y)$ by
$-i$ and $h_q(X,J_qY)$ by $i$. The two forms cannot agree
identically on a nonzero positive subspace.
\end{proof}

This gives an exact criterion for transporting the local operator
construction to a native pairing. Let $V$ and $V^!$ be complex
vector spaces, let $B:V\times V^!\to\mathbb C$ be bilinear,
let $h$ be a Hermitian form on $V$ linear in its first variable,
and let $c:V\to V^!$ be conjugate-linear. A realization of these
data in the GNS carrier consists of complex-linear maps $f,g$
with
\begin{equation}
 B(v,a)=B_q(fv,ga),\qquad
 h(v,w)=h_q(fv,fw),\qquad gc=J_qf.
 \label{v4-arith:thermal:pairing-realization-data}
\end{equation}
The first two conditions concern different pairings on specified
carriers; the third concerns the anti-linear comparison map.
They are checked independently before composing them.

\begin{proposition}[pairing transfer]
\label{v4-arith:thermal:pairing-transfer}
If \eqref{v4-arith:thermal:pairing-realization-data} holds, then
$B(v,cw)=h(v,w)$. If $f$ is injective, the transported Hermitian
form is positive definite. The displayed conditions do not follow
from equality of partition functions.
\end{proposition}

\begin{proof}
The three specified maps give
\[
 B(v,cw)=B_q(fv,gcw)=B_q(fv,J_qfw)=h_q(fv,fw)=h(v,w).
\]
If $f$ is injective, $h(v,v)=\|fv\|_2^2$ vanishes only for
$v=0$. The phase obstruction and the vacuum commutator show that
a scalar character omits operations used in these equations;
no map in \eqref{v4-arith:thermal:pairing-realization-data} is
constructed by a scalar equality alone.
\end{proof}

\section{The global trace and the native categorical comparison}

For all primes at once, the occupation basis is indexed by positive
integers through $\mathbf n\mapsto\prod_p p^{n_p}$, with finite
support. Its Hilbert completion is $\ell^2(\mathbb N_{\geq1})$ and
the energy is $\log n$ on the vector indexed by $n$. The global
thermal density is nuclear exactly for real $\beta>1$, because
its trace is $\sum_{n\geq1}n^{-\beta}$. The convergence for
$\beta>1$ and divergence for $0<\beta\leq1$ follow by comparison
with the integral of $x^{-\beta}$ over $[1,\infty)$.
Unique factorization gives its Euler product. The compatible
normalized states on the algebraic union of the tensor shift
algebras $\mathcal T_F$ still exist for every $\beta>0$: positivity and
normalization are checked in a finite tensor factor containing
the given element. Existence of this local-state system is not
the assertion of a nuclear global Gibbs density at those temperatures.

The unnormalized traces behave differently under a prime-set
inclusion: tensoring an observable with the identity multiplies
its thermal trace by the partition factor of each new prime.
The coefficient maps $m\mapsto m\rho_{F,\beta}$ are likewise not
intertwined by the ordinary inclusions $K\mapsto K\otimes1$.
They are intertwined by $K\mapsto K\otimes\rho_{G\setminus F,\beta}$,
which is a coefficient-module construction and is not an algebra
inclusion. These distinctions must persist in any global comparison.

The native arithmetic realization asks for more than these local
operator and finite Hecke constructions. Its proposed source is
the Hochschild homology of a restricted tensor product of local
Iwahori convolution categories, with Frobenius data. A realization
must first define that tensor product and its trace in a specified
coefficient category, then construct maps from the local data
used here into it. A comparison with BPS cohomology further needs
actual quivers and polynomial potentials, cohomological and
dimension gradings, compatible prime-set maps, and an intertwining
of the trace and energy operations.

The local results identify three conditions on such a construction.
The represented prime phase quotients do not admit the naive
phase-preserving prime-set inclusions; the two-prime calculation
requires a larger compatible source before a global trace can
be compared.
An ordinary cyclic functional on the full prime shift algebra
cannot equal the finite-temperature Gibbs functional; the
imaginary-time twist or the nuclear coefficient module is necessary
for the comparison constructed here. A multiplication-preserving
map from the finite Iwahori Hecke algebra to the connected
polynomial Jordan target cannot preserve its trace; its idempotent
sectors cannot be removed. A different categorical or critical
Hall target is not excluded by these arguments. Constructing that
target and its native arithmetic comparison remains open.

